\documentclass[11pt]{article}
\usepackage[letterpaper,margin=1in]{geometry}
\usepackage{fontspec}
\usepackage{amsmath,amssymb}
\usepackage{graphicx,booktabs,xcolor,microtype}
\usepackage[numbers,sort&compress]{natbib}
\usepackage[font=small,labelfont=bf]{caption}
\usepackage[hidelinks,pdfusetitle]{hyperref}
\usepackage{enumitem}
\setlist{itemsep=1pt,topsep=3pt}
\graphicspath{{figures/}}
\newcommand{\ShuffleSnap}{ShuffleSnap}
\newcommand{\panel}[2]{\par\begin{center}\begin{minipage}{\linewidth}\centering
  \includegraphics[width=\linewidth]{#1}\captionof{figure}{#2}\end{minipage}\end{center}}

\title{\ShuffleSnap: strided windowed exact reassignment\\for point-to-grid layout}
\author{Kyle McDonald}
\date{September 2026}

\begin{document}
\maketitle

\begin{abstract}
Giving each point of a 2D embedding its own grid cell with the least total squared
displacement is a linear assignment problem too large for dense solvers at $10^5$ points.
\ShuffleSnap, inspired by Shellsort, solves small exact assignment problems in windows whose
cells lie a stride apart, and shrinks the stride to one. Against certified optima it comes
within hundredths of a percent on clustered layouts, including real UMAP embeddings, in
seconds, closer than BSP-OT and RasterFairy on most inputs. Its errors are loops of one-cell shifts that
no window can see, which make uniform points, concave masks and surplus cells hard.
\end{abstract}

\section{Introduction}

An image collection is often shown by embedding the items in the plane, for example with
t-SNE or UMAP~\citep{vandermaaten2008tsne,mcinnes2018umap}, and placing each thumbnail in
its own grid cell. Rounding every point to its nearest cell causes collisions and holes, so
the cells must be assigned jointly: we seek the assignment $\sigma$ of points $p_i$ to
distinct cells $c_{\sigma(i)}$ that minimizes $C(\sigma)=\sum_i \lVert p_i - c_{\sigma(i)}\rVert^2$.
This is a linear assignment problem, and a discrete optimal-transport
problem~\citep{peyre2019computational}. Dense exact
solvers~\citep{kuhn1955hungarian,jonker1987shortest} need the full cost matrix, 80\,GB at
$10^5$ points. Practical tools instead split the cloud recursively, like
RasterFairy~\citep{klingemann2015rasterfairy}, or merge randomized partitions, like
BSP-OT~\citep{genest2025bspot}.

\ShuffleSnap\ borrows from Shellsort~\citep{shell1959highspeed}, which sorts interleaved
subsequences at a shrinking gap. It solves only small assignment problems, but exactly,
with their cells spread across the grid first and brought together later
(Figure~\ref{fig:schematic}). We treat the idea as a hypothesis, and measure its true gaps
with a solver that proves optimality on the same instances.

\panel{fig_schematic}{(a)~A window collects the points that occupy its cells and
reassigns them among the same cells by an exact linear assignment. (b)~At stride $s$ a
window takes every $s$-th row and column: it keeps 64 cells but spans $7s+1$ cells per
axis, so points can travel far. (c)~Each subgrid is tiled four times, shifted by half a
window, so that points can cross every window border.\label{fig:schematic}}

\clearpage

\section{The method}

\ShuffleSnap\ improves a valid assignment in place:
\begin{enumerate}
\item \textbf{Start} from a seeded random permutation of the points over the cells.
\item \textbf{Stride.} At stride $s$, cells whose row and column have the same residues
  modulo $s$ form one of $s^2$ interleaved subgrids. Each subgrid is cut into windows of
  $8\times8$ of its cells, which span $7s+1$ cells of the grid per axis.
\item \textbf{Solve.} In every window, the current occupants are reassigned among the
  window's cells, empty ones included, by an exact linear assignment on squared distances
  to the original coordinates. The result replaces the old assignment only if it is cheaper
  by more than a relative $10^{-12}$.
\item \textbf{Shift.} A round tiles every subgrid four times, shifted by $(0,0)$,
  $(0,4)$, $(4,0)$ and $(4,4)$. Windows that cross an edge are clipped, never dropped.
\item \textbf{Schedule.} On an axis of length $L$ the first stride is the smallest power
  of two with $8s\ge L$, so that the first windows span the grid. After each round the
  stride is divided by $\sqrt2$ and rounded; at stride 1, rounds repeat until one moves no
  point.
\end{enumerate}
Because the old assignment is feasible for every window, the total cost never increases,
the layout is valid after every step, and ties cannot cause cycling. The windows of one
tiling are disjoint, so they run in parallel with a result that does not depend on the
number of threads. Each level costs about the same, linear in the number of cells. The
final no-move round is a fixed point of the windows, not a proof of optimality.

Rectangular grids use one stride per axis, masked cells never enter a window, and empty
cells do, so with more cells than points the solver also chooses which cells stay empty.
The kernel is C++ with OpenMP and a Python interface, needs about 55 bytes per point, and
skips windows whose occupants have not changed. It solves each window exactly from starting
duals derived in closed form from per-axis optimal transport, so that the points of a distant
cluster do not all compete for one corner cell. A 70-line serial version in NumPy and SciPy
reproduces its assignments exactly. We chose the window side and the stride ratio on
development instances. Larger windows were better in every test but slower: each step from
$8\times8$ to $12\times12$ to $16\times16$ cost 2--2.6 times the time and removed one half to
two thirds of the remaining excess. We offer \emph{fast} ($5\times5$), \emph{balanced}
($8\times8$, default) and \emph{quality} ($12\times12$); fast and balanced take about as long
as BSP-OT with 16 and 64 plans, so the two compare at equal time. Dividing the stride by $\sqrt2$ halves the area a window spans from
one level to the next, so each level starts from a layout already sorted at nearly its own
scale; it was the best or close to the best ratio at every window size.

\panel{fig_progress}{Progress on 10{,}000 points of a UMAP embedding of MNIST
digits~\citep{lecun1998gradient}. Every point keeps the color of its input position, and
the dark bands are graph-paper lines drawn on the input: both travel with the points
(every second stride shown). The coarse strides sort the layout globally, stride 1 polishes
it. Numbers are the mean squared
displacement in cells$^2$, next to the certified optimum.\label{fig:progress}}

\clearpage

\section{Certified optima}

Gaps need a true optimum, but at $10^5$ points no dense solver fits in memory. For any cell potentials $v$, however,
\[
  \mathrm{LB}(v)=\sum_{i=1}^N \min_j\bigl(\lVert p_i-c_j\rVert^2-v_j\bigr)+\sum_{j=1}^M v_j-(M-N)\max_j v_j
\]
is a lower bound on the optimal cost, by weak duality of the
assignment linear program, where $M\ge N$ is the number of cells; it is tight for optimal
potentials. Our solver therefore works on a sparse problem and checks it against the dense
one. It solves a min-cost flow on candidate edges (each point's nearest cells and the
neighbourhood of its cell in a starting layout) with OR-Tools~\citep{ortools}, computes
potentials from that solution by shortest paths, evaluates $\mathrm{LB}$ over \emph{all}
cells, adds every edge that the bound shows to be missing, and repeats until the bound
meets the cost. The result is then the exact optimum of the dense problem; the start affects
only the time. This is the core-and-price idea of sparse
assignment and transport solvers~\citep{volgenant1996linear,schmitzer2016sparse}. Wherever
SciPy's dense solver fits, both found the same optimum: on 600 instances of up to 2{,}500
points and on all ten distributions at $10^4$ points. Started from \ShuffleSnap's layout, it certified all 202
benchmark instances up to $2\times10^5$ points, and uniform instances up to $10^6$.

\section{Setup}

All runs used one workstation (Intel i9-12900KS, 24 hardware threads). Single-thread runs
were pinned to one performance core; every run was a fresh process, and times include each
method's own preparation but not generating the instance. Design decisions used one set of
instances, and all results below use three held-out ones. The inputs are ten synthetic
distributions with $10^5$ points on near-square grids (Table~\ref{tab:main}), the complete
UMAP embeddings of MNIST and Fashion-MNIST~\citep{xiao2017fashion} with
70{,}000 points each, and masks. Every method received identical coordinates and cells. We
compare BSP-OT through its official Python binding with 16 and 64 merged plans, and
RasterFairy in its current upstream version, which solves leaf blocks exactly and is far
better than its 2019 PyPI release. Table~\ref{tab:main} lists the results for all ten distributions.

\begin{center}
\begin{minipage}{\linewidth}\centering
\captionof{table}{Excess over the certified optimum (\%) at $N=100{,}000$, mean of three
held-out instances, one thread. The last column is the certified solver's time, started
from \ShuffleSnap; the bottom row is the median time of each method.\label{tab:main}}
\small
\input{tables/tab_main_compact_100k_t1}
\end{minipage}
\end{center}

\clearpage

\section{Quality, time and scale}

\panel{fig_summary}{Excess over the certified optimum against time on one core. Points
are means over three instances (for MNIST, three runs on the one embedding); bars span
their range. The three \ShuffleSnap\ presets use $5\times5$, $8\times8$ (default) and
$12\times12$ windows from left to right. The vertical line is the time of the certified exact
solver.\label{fig:summary}}

On clustered data, the kind of layout that embeddings produce, \ShuffleSnap\ is close to
optimal (Figure~\ref{fig:summary}, Table~\ref{tab:main}). At $10^5$ points the default
preset is within 0.06\% of the optimum on nine of the ten distributions, and within
0.005\% on six, in a median 2.3\,s on one core. At equal time it beats BSP-OT: the fast preset
(0.77\,s) is closer to the optimum than BSP-OT with 16 plans (0.71\,s) on eight distributions,
and the default than 64 plans (2.9\,s) on nine; on uniform points they reach 2.6\% and 0.84\%,
against 45\% and 32\%. BSP-OT is ahead only on the near-line and the 90/10 clusters.
RasterFairy is 17--425\% above the optimum on seven distributions. Only below 0.12--0.36\,s on
one core, before \ShuffleSnap\ has finished its coarse levels, is BSP-OT's result the better
one.

With 24 threads \ShuffleSnap\ runs 7--14 times faster and returns bit-identical layouts:
$10^5$ points take 0.09\,s with the fast preset and 0.19\,s with the default, less than
BSP-OT's 0.13\,s and 0.41\,s. Time grows nearly linearly with $N$; $10^6$ uniform points take
21\,s on one core and 1.5\,s on 24 threads with the default, using about 55 bytes per point, against about 250 and 670 for BSP-OT and
RasterFairy. A billion points converged in 1\,h\,43\,min using 53\,GiB, passing a six-hour run of an
earlier implementation after 44 minutes; a single-plan BSP-OT ended 8.2\% above. The exact solver scales worse but further than we expected. Seeded with
\ShuffleSnap's answer, it proved optimality at $10^5$ points in 17\,s to 23\,min on one
core, with a median near two minutes, and at $10^6$ uniform points in 11\,min.

What mattered in the design: larger windows are the main quality knob, the coarse levels
are essential, without shifted tilings the method fails outright, and the $\sqrt2$ schedule
is best or close to best among the stride ratios we tried. Structured starts, extra
polishing at stride 2, merging independent runs and a Jonker--Volgenant window solver did
not pay off. Starting each window solve from closed-form transport duals made it 2--5 times
faster, which buys larger windows at equal time.

\paragraph{Masks and surplus cells.} With 20{,}000 points on disks and rings, the default
preset stays within 0.07\% of the optimum (fast: 0.14\%), where BSP-OT is 2.5--9.3\% and
RasterFairy 11--54\% above it. Surplus cells, where a method must also choose which cells
stay empty, hurt the comparators most: with 10\% surplus, BSP-OT's partial matching is
18--109\% above the optimum and RasterFairy, which removes surplus cells at random,
72--350\%, while the default stays within 0.91\%. Concave masks and fewer surplus cells are
harder (Section~\ref{sec:fail}).

\clearpage

\section{Real embeddings: displacement is not neighbourhood}

\panel{fig_gallery_mnist}{All 70{,}000 MNIST digits on a $265\times265$ grid. Colors and
graph-paper lines follow each point's input position; numbers are the mean squared
displacement in cells$^2$. The bottom row shows lattice precision, the fraction of a point's
occupied neighbour cells that hold one of its 8 nearest input neighbours.\label{fig:gallery}}

On both embeddings the default preset is within 0.002\% of the certified optimum, and its
layout is visually indistinguishable from the optimal one: the graph-paper lines bend in
the same places (Figure~\ref{fig:gallery}). The two comparators fail in opposite ways.
BSP-OT is within 2.4\% (MNIST) and 6.7\% (Fashion-MNIST) of the optimal displacement, yet it
halves both neighbourhood measures (the fraction of each point's 10 nearest input neighbours
among its 10 nearest grid neighbours is 0.31--0.33, against 0.61--0.65 for the optimum): its merged partitions follow the optimal map on average
but scramble points locally, which frays the lines. RasterFairy moves points 20--25\% more
than necessary, but keeps local order inside its slices and breaks it only along their
seams, so its neighbourhood measures stay close to the optimum's. Even the optimum tears
the layout along cluster boundaries, visible as pale lines in the bottom row. Minimizing
displacement and preserving neighbourhoods are different goals; neither implies the other.

\clearpage

\section{Where it fails}
\label{sec:fail}

\panel{fig_failures_row}{(a,\,b)~Where \ShuffleSnap\ and the optimum disagree on uniform
points, the difference is a few long closed loops of one-cell moves (loops longer than 30
points in blue), and at $10^5$ points loops of more than 300 points carry most of the excess.
(c,\,d)~Points drawn uniformly inside a blocky S-shaped mask, whose optimum moves every point
by less than a cell: $6\times6$ windows with halving strides route points across the concave
bays, where they line up along the edges; the default preset leaves fewer such lines. Numbers are mean squared
displacements.\label{fig:failures}}

\ShuffleSnap\ is a local search, and its failures follow from that. The excess sits in long
loops in which every point is shifted by about one cell (Figure~\ref{fig:failures}a,b).
Such a loop is optimal inside every window that sees only part of it, and a strided window
cannot shift a point by one cell. We built inputs on which the method never moves while the
optimum costs zero, so its relative excess is unbounded. Uniform points are the hardest
common input: the default's excess grows with $N$, from 0.1\% at $10^3$ points to 0.84\% at
$10^5$ and 1.3\% at $10^6$. On masks with concave bays, coarse schedules route points across
the bays (Figure~\ref{fig:failures}c,d). The default reduces this, and with uniform points on
the S mask it is 1.4\% above the optimum (BSP-OT: 9.3\%). A mask can even cut two usable cells off from each other, when no
scheduled window contains both. Grids with a few surplus cells are hard too, because the
solver must also choose where the holes go (3.1\% with less than a row of them).

\paragraph{When to use it.} For interactive layouts of $10^5$ points, the default preset
gives near-optimal results in about 0.2\,s on a multicore machine (fast: 0.09\,s). When a proven optimum
is worth minutes, the certified solver started from \ShuffleSnap\ is the better choice.
Below 0.12--0.36\,s on one core, BSP-OT gives the best available result. If the position
of empty cells does not matter, an exact-fit mask is easier than surplus cells for every
method. The code, raw results and figure scripts are released in two repositories.

\paragraph{Limits of this study.} All measurements come from one workstation, CPU only;
BSP-OT also has a GPU implementation that we did not test. The objective is squared
Euclidean displacement in 2D, the real data are two UMAP embeddings with one parameter
setting, and each comparator ran at representative settings, not at every setting.
\label{page:body-end}

\clearpage
\bibliographystyle{plainnat}
\bibliography{refs}
\end{document}
